\chapter{Методологија за тестирање и резултати}
\label{chap:results}

\section{Доказна скала}

За да не се изедначат source code, тест и физичко однесување, се користи
следната доказна скала:

\begin{enumerate}
  \item \status{имплементирано} --- source/configuration постои и е
  прегледана;
  \item \status{unit/offline тестирано} --- deterministic тестот поминал без
  hardware;
  \item \status{потврдено во Gazebo} --- интегрираниот simulator graph покажал
  мерливо и, кога е наведено, визуелно однесување;
  \item \status{физички извршено} --- exact runner/profile бил извршен на
  роботот со датирана телеметрија;
  \item \status{визуелно прифатено} --- операторот го оценил движењето како
  соодветна емоција;
  \item \status{deployed/offline-tested, live-unvalidated} --- binary/source е
  поставен и тестовите поминале, но новата behavior branch не е повторно
  физички активирана; и
  \item \status{планирано/неимплементирано}.
\end{enumerate}

Topic publication, wire packet, sent joint target или joint movement не се
автоматски доказ за правилна емоционална реакција. Потребни се корелирани
intent, transport, state, joints/IMU/contact, release и operator observation
соодветни на тврдењето.

\section{Симулациска методологија}

Симулациската проверка се извршува во чист, единствен ROS/Gazebo graph. По
Release build се стартува headless integrated launch, се чекаат controller
readiness и prepared stance и се инјектираат exact \code{event:<emotion>}
влезови. Се собираат:

\begin{itemize}
  \item conversation event и emotion-state со turn/sequence correlation;
  \item expression command и generation-aware action;
  \item safe command и final Joy transport;
  \item model position, torso height/tilt и сите 12 joint positions;
  \item foot contact topics, controller states и safety status; и
  \item final zero command, residual speed и teardown markers.
\end{itemize}

Првиот sweep проверува entrance, rapid changes, manual arbitration, watchdog и
no-hardware boundary. Вториот го задржува секој профил низ два idle cycles и
проверува movement, 0.10 m containment, 0.05 rad joint margin, 0.30 rad HipX
anti-splay, mirrored stomp, exact-neutral holds и отсуство на NaN/fall/error.
Live visual claim бара блиска Gazebo камера покрај telemetry; headless бројките
сами не докажуваат читливост.

\section{Физичка методологија}

Физичкиот preflight е read-only до моментот на одделна operator authorization.
Се проверуваат model/software, exact initial state 1, gait/motion 0/0, zero
errors, battery floor, attitude, fresh \code{0x0901}/\code{0x0906}, centered
Retroid axes, STOP false, active passive services и отсуство на owner/runner.
Потоа official runner извршува \cref{eq:official-sequence} и задржува единствен
lease.

За planted profile се проверуваат joint bounds, tracking, attitude, state,
feedback и final release. За lifted paw дополнително се воспоставува
distinct-tick baseline, се проверуваат unload, non-target support, landing и
four-foot relatch. Miss не се „поправа“ со повисока амплитуда: paw се спушта,
се прави recovery, а исходот се заведува како failure или строго дефиниран
recoverable contact miss.

Визуелната евалуација е намерно одвоена. Операторот оценува дали motion reads as
Neutral/Joy/Sadness/Fear/Anger. Таа оценка не може да поништи safety gate, но е
неопходна за профилот да се нарече емоционално прифатен.

\section{Evidence matrix}

\Cref{tab:evidence-matrix} ги сумира најважните резултати. Сите project-row
резултати се од 20 септември 2026 година и од commit-ите наведени во изјавата
за опфат, освен каде што самиот датиран record јасно означува претходна
working-tree основа. Во таков случај record-от е историски доказ, а референтниот
snapshot ја содржи финалната документирана состојба.

\begingroup
\scriptsize
\setlength{\tabcolsep}{2pt}
\setlength{\LTleft}{0pt}
\setlength{\LTright}{0pt}
\begin{longtable}{@{}L{0.18\textwidth}L{0.11\textwidth}L{0.18\textwidth}L{0.29\textwidth}L{0.18\textwidth}@{}}
\caption{Evidence matrix за референтниот snapshot}\label{tab:evidence-matrix}\\
\toprule
\textbf{Компонента} & \textbf{Тип} & \textbf{Тест/команда} & \textbf{Средина и резултат} & \textbf{Датум/ограничување} \\
\midrule
\endfirsthead
\toprule
\textbf{Компонента} & \textbf{Тип} & \textbf{Тест/команда} & \textbf{Средина и резултат} & \textbf{Датум/ограничување} \\
\midrule
\endhead
\bottomrule
\endfoot
Release catkin build & офлајн & \code{emotion-clean-build} & Ubuntu 20.04/Noetic Docker; поминат со постојни upstream warnings & 2026-09-20; Lite3\_VMC \code{785a384} \\
Python unit suites & офлајн & \code{emotion-unit-tests} & Noetic container; 58 тестови поминати & 2026-09-20; без physical claim \\
C++ trajectories & офлајн & controller test binary & Noetic container; 5 тестови поминати & 2026-09-20; simulation mechanics \\
OpenAI process boundary & офлајн & \code{emotion-openai-}\allowbreak\code{bridge-test} & Python 3.12 sidecar + Noetic client; 1 тест, без key/network & 2026-09-20 \\
ROS contracts/ordering & офлајн ROS & \code{emotion-ros-tests} & ROS master/container; 1 suite, 0 failures/errors & 2026-09-20 \\
Gazebo profiles & симулација & два headless sweeps & Gazebo/13 controllers; сите 9 profiles и 12 joints & 2026-09-20; не е hardware \\
No-hardware boundary & инспекција и тест & Gazebo E2E & Docker network/process scan; без real executable/robot UDP & 2026-09-20 \\
Hardware offline gate & офлајн & \code{verify-hardware-}\allowbreak\code{offline} & Hardware catkin container; 9 native suites, 29 Python tests, build/integration pass & 2026-09-20; не е commissioning \\
aarch64 runner & build/native & clean ARM build и suites & Perception computer, official SDK; 9 suites pass, binaries installed & 2026-09-20; installation не е motion \\
Neutral breathing & физичко и визуелно & continuous runner & Physical Lite3; 27 cycles at handoff, прифатено & 2026-09-20; exact tested setup \\
Joy profile & физичко и визуелно & 5 s paw suite + chat loops & Physical Lite3; unload/landing на двете paws, прифатено & 2026-09-20; latest recovery не е live revalidated \\
Sadness profile & физичко и визуелно & single + 2-cycle bow & Physical Lite3; four supports, прифатено & 2026-09-20; paw variant withdrawn \\
Fear profile & физичко и визуелно & 3 planted cycles & Physical Lite3; 15 s, clean release, прифатено & 2026-09-20; paw variant rejected \\
Anger profile & физичко и визуелно & 3-cycle reset suite & Physical Lite3; 6 placements, four-foot landings & 2026-09-20; latest recovery не е live revalidated \\
Exact-neutral retarget & физичко & normal chat transitions & Physical Lite3; accepted/fallback retargets pass & 2026-09-20; consolidated verdict pending \\
Battery interlock & физичко fail-safe & 75\% session floor & Physical Lite3; release на 74\%, без owner leak & 2026-09-20; намерен abort \\
Joy persistent recovery & deploy и offline & full gate + ARM deploy & Container + perception computer; pass, SHA prefix \code{aa449b} & 2026-09-20; без последователно motion \\
Anger persistent recovery & deploy и offline & full gate + ARM deploy & Container + perception computer; pass, SHA prefix \code{e72a2d} & 2026-09-20; без последователно motion \\
\end{longtable}
\endgroup


\section{Толкување на резултатите}

Резултатите поддржуваат четири заклучоци. Прво, архитектурното разделување е
функционално: истиот EmotionEngine и state contract ги хранат Gazebo и
hardware resolver-от, без simulator command да стане hardware command. Второ,
сите девет категории се изразливи и автоматски проверливи во Gazebo. Трето,
пет категории имаат физички и визуелно прифатени реакции, но целата normal
configuration намерно останува Neutral-only. Четврто, contact gates и
transition contract не се формалност: токму тие откриле support redistribution,
unload misses, unsafe takeover и attitude/battery boundaries.

Постојат и важни негативни резултати. High-level height path не докажал
значајно periodic motion; првите Joy/Sadness/Fear candidates не ја пренесувале
саканата емоција; Anger endurance и recovery policies откриле реални пропусти.
Вклучувањето на овие резултати ја подобрува репродуцибилноста и спречува
погрешна генерализација.
